\chapter{भिन्नें: तुलना और जोड़}\label{ch:g7:fractions}

पिछले साल भिन्नें हिस्सों के रूप में और ठीक भागफलों के रूप में
आई थीं (\cref{ch:g6:fractions})। यह अध्याय उनके साथ \emph{गणना}
करना सिखाता है: बराबर भिन्नें पहचानना, उनकी तुलना करना, उन्हें
जोड़ना और घटाना --- फ़िलहाल आसान हरों के साथ; आम स्थिति
\cref{ch:g9:fractions} में आएगी।

\section{बराबर भिन्नें}

\begin{theorem}[बराबर भिन्नें]\label{thm:g7:fractions:equal}
\omterm{def:g6:fractions:def}{भिन्न} का \omterm{def:g4:fractions:def}{अंश} और हर दोनों एक ही शून्येतर संख्या से गुणा कर दिए
जाएँ, या भाग दे दिए जाएँ, तो \omterm{def:g6:fractions:def}{भिन्न} नहीं बदलती:
\[
\frac{a}{b} = \frac{a \times k}{b \times k} .
\]
\end{theorem}
\admitted

\begin{example}\label{ex:g7:fractions:equal}
$\dfrac{3}{4} = \dfrac{3 \times 5}{4 \times 5} = \dfrac{15}{20}$, और
दूसरी दिशा में $\dfrac{42}{30} = \dfrac{42 \div 6}{30 \div 6}
= \dfrac{7}{5}$ ($6$ से सरल की हुई)।

यह जाँचने के लिए कि $\dfrac{8}{12}$ और $\dfrac{10}{15}$ बराबर हैं
या नहीं, दोनों को सरल करो: $\dfrac{8}{12} = \dfrac{2}{3}$ और
$\dfrac{10}{15} = \dfrac{2}{3}$ --- हाँ, बराबर।
\end{example}

\section{भिन्नों की तुलना}

\begin{method}[दो भिन्नों की तुलना]\label{met:g7:fractions:compare}
\begin{enumerate}
  \item \emph{एक ही हर}: अंशों की तुलना करो:
        $\frac{5}{9} < \frac{7}{9}$।
  \item \emph{एक हर दूसरे का \omterm{def:g5:division:multiple}{गुणज} है}: बड़े हिस्सों वाली \omterm{def:g6:fractions:def}{भिन्न} को
        दूसरे हर के साथ फिर से लिखो, फिर तुलना करो:
        $\frac{5}{6}$ बनाम $\frac{7}{12}$: $\frac56 = \frac{10}{12} >
        \frac{7}{12}$।
  \item जहाँ हो सके \emph{$1$ या $\frac12$ से तुलना करो}:
        $\frac{9}{8} > 1 > \frac{7}{9}$ से $\frac98$ और $\frac79$ का
        मामला पल भर में तय हो जाता है।
\end{enumerate}
\end{method}

\begin{omfigure}
\begin{tikzpicture}[scale=0.62]
  \foreach \i in {0,...,5} \draw[omExo] (\i*2,0) rectangle (\i*2+2,1.1);
  \foreach \i in {0,...,4} \fill[omDef!35] (\i*2+0.05,0.05)
    rectangle (\i*2+1.95,1.05);
  \draw[thick] (0,0) rectangle (12,1.1);
  \node[right, font=\small] at (12.3,0.55) {$\frac56$};
  \begin{scope}[yshift=-1.9cm]
  \foreach \i in {0,...,11} \draw[omExo] (\i,0) rectangle (\i+1,1.1);
  \foreach \i in {0,...,6} \fill[omThm!35] (\i+0.05,0.05)
    rectangle (\i+0.95,1.05);
  \draw[thick] (0,0) rectangle (12,1.1);
  \node[right, font=\small] at (12.3,0.55) {$\frac{7}{12}$};
  \end{scope}
\end{tikzpicture}

{\small जुड़वाँ पट्टियों पर $\frac56$ और $\frac{7}{12}$ की तुलना:
छठे हिस्सों को आधा-आधा काटने पर दिखता है कि $\frac56 = \frac{10}{12}$,
जो $\frac{7}{12}$ से बड़ा है।}
\end{omfigure}

\section{जोड़ना और घटाना}

\begin{theorem}[एक ही हर]\label{thm:g7:fractions:add}
एक ही हर वाली भिन्नों को जोड़ने (या घटाने) के लिए उनके \omterm{def:g4:fractions:def}{अंश} जोड़
(या घटा) दिए जाते हैं:
\[
\frac{a}{d} + \frac{b}{d} = \frac{a + b}{d},
\qquad
\frac{a}{d} - \frac{b}{d} = \frac{a - b}{d} .
\]
\end{theorem}

\begin{proof}[क्यों]
$\frac1d$ आकार के $a$ हिस्से और उतने ही आकार के $b$ हिस्से मिलकर
उसी आकार के $a + b$ हिस्से बनाते हैं।
\end{proof}

\begin{method}[अलग-अलग हर]\label{met:g7:fractions:add}
जब एक हर दूसरे का \omterm{def:g5:division:multiple}{गुणज} हो:
\begin{enumerate}
  \item छोटे हर वाली \omterm{def:g6:fractions:def}{भिन्न} को फिर से लिखो ताकि दोनों का हर बड़ा
        वाला हो जाए (\cref{thm:g7:fractions:equal});
  \item अंशों को जोड़ो या घटाओ;
  \item हो सके तो उत्तर को सरल करो।
\end{enumerate}
\end{method}

\begin{example}\label{ex:g7:fractions:add}
$\dfrac{3}{4} + \dfrac{5}{8}$ की गणना, क़दम दर क़दम:
\begin{enumerate}
  \item $8$, $4$ का \omterm{def:g5:division:multiple}{गुणज} है: $\dfrac34 = \dfrac{3 \times 2}{4 \times 2}
        = \dfrac{6}{8}$ लिखो;
  \item जोड़ो: $\dfrac{6}{8} + \dfrac{5}{8} = \dfrac{11}{8}$;
  \item सरल करने को कुछ नहीं: उत्तर $\dfrac{11}{8}$ है, यानी
        $1 + \dfrac38$।
\end{enumerate}
एक और: $2 - \dfrac{4}{5} = \dfrac{10}{5} - \dfrac{4}{5} =
\dfrac{6}{5}$ (पहले पूर्ण संख्या को हर $5$ के ऊपर लिख लो)।
\end{example}

\begin{example}[एक जानी-पहचानी ग़लती]\label{ex:g7:fractions:trap}
$\dfrac12 + \dfrac13$, $\dfrac{2}{5}$ \emph{नहीं} है! अंशों और हरों
को अलग-अलग जोड़ना ग़लत है --- उत्तर $\frac25$ तो $\frac12$ से भी
छोटा है, जो बेतुका है जब $\frac12$ में कुछ धनात्मक \emph{जोड़ा} जा
रहा हो। (सही \omterm{def:g1:addition:def}{योगफल} $\frac56$ है, जिसके लिए उभयनिष्ठ हर $6$ चाहिए;
आम तरीक़ा \cref{ch:g9:fractions} में है।)
\end{example}

\section{भिन्न को किसी संख्या से गुणा करना}

\begin{proposition}[भिन्न गुना संख्या]\label{prop:g7:fractions:mult}
किसी संख्या $k$ के लिए:
\[
k \times \frac{a}{b} = \frac{k \times a}{b} .
\]
किसी राशि का $\frac ab$ \emph{भाग} लेने का मतलब है उसे $\frac ab$
से गुणा करना।
\end{proposition}
\admitted

\begin{example}\label{ex:g7:fractions:mult}
$6 \times \dfrac{2}{3} = \dfrac{12}{3} = 4$: दो-तिहाई का छह गुना
चार पूरे होते हैं। और ``$60$ यूरो का $\frac34$'' है
$\frac34 \times 60 = \frac{180}{4} = 45$ यूरो --- वही उत्तर जो $4$
से भाग देकर फिर $3$ से गुणा करने पर मिलता है
(\cref{met:g6:fractions:ofquantity})।
\end{example}

\section{अभ्यास}

\begin{exercise}[$\star$]\label{exo:g7:fractions:1}
पूरा करो: $\dfrac{2}{5} = \dfrac{?}{20}$; \quad
$\dfrac{18}{24} = \dfrac{3}{?}$; \quad
$\dfrac{7}{3} = \dfrac{28}{?}$; \quad
$\dfrac{?}{9} = \dfrac{20}{36}$.
\end{exercise}

\begin{exercise}[$\star$]\label{exo:g7:fractions:2}
जितना हो सके सरल करो:
$\dfrac{12}{16}$; \quad $\dfrac{30}{45}$; \quad $\dfrac{27}{36}$;
\quad $\dfrac{48}{12}$.
\end{exercise}

\begin{exercise}[$\star$]\label{exo:g7:fractions:3}
क्या $\dfrac{15}{35}$ और $\dfrac{21}{49}$ बराबर हैं? और
$\dfrac{16}{28}$ तथा $\dfrac{20}{36}$? सरल करके कारण बताओ।
\end{exercise}

\begin{exercise}[$\star$]\label{exo:g7:fractions:4}
तुलना करो (सिर्फ़ उत्तर नहीं, तर्क भी लिखो):
\[
\frac{7}{11} \text{ और } \frac{9}{11}; \qquad
\frac{5}{6} \text{ और } \frac{11}{18}; \qquad
\frac{13}{12} \text{ और } \frac{19}{20}.
\]
\end{exercise}

\begin{exercise}[$\star$]\label{exo:g7:fractions:5}
हो सके तो सरल करते हुए गणना करो:
\[
\frac{5}{9} + \frac{2}{9}, \qquad
\frac{11}{7} - \frac{4}{7}, \qquad
\frac{5}{8} + \frac{7}{8} .
\]
\end{exercise}

\begin{exercise}[$\star$]\label{exo:g7:fractions:6}
\cref{met:g7:fractions:add} की मदद से गणना करो:
\[
\frac{2}{3} + \frac{5}{6}, \qquad
\frac{7}{10} - \frac{2}{5}, \qquad
\frac{3}{4} + \frac{5}{12} .
\]
\end{exercise}

\begin{exercise}[$\star$]\label{exo:g7:fractions:7}
गणना करो:
\[
3 - \frac{5}{4}, \qquad
1 + \frac{3}{8}, \qquad
2 - \frac{7}{6} .
\]
\end{exercise}

\begin{exercise}[$\star$]\label{exo:g7:fractions:8}
गणना करो:
\[
5 \times \frac{3}{10}, \qquad
\frac{7}{8} \times 4, \qquad
45 \text{ का } \frac{2}{3} .
\]
\end{exercise}

\begin{exercise}[$\star\star$]\label{exo:g7:fractions:9}
मार्क ने एक पिज़्ज़ा का $\frac14$ खाया और जूली ने उसी पिज़्ज़ा का
$\frac38$। दोनों ने मिलकर पिज़्ज़ा का कितना भाग खाया? कितना भाग
बचा?
\end{exercise}

\begin{exercise}[$\star\star$]\label{exo:g7:fractions:10}
एक बोतल में $\frac34$ L रस है। लेआ दो बार $\frac16$ L उँडेलती है।
कितना रस बचा? (उभयनिष्ठ हर: $12$।)
\end{exercise}

\begin{exercise}[$\star\star$]\label{exo:g7:fractions:11}
\cref{ex:g7:fractions:trap} वाले तर्क से समझाओ कि सही \omterm{def:g1:addition:def}{योगफल} निकाले
बिना ही क्यों $\frac35 + \frac12$, $\frac{4}{7}$ नहीं हो सकता।
\end{exercise}

\begin{exercise}[$\star\star\star$]\label{exo:g7:fractions:12}
\omterm{def:g1:addition:def}{योगफल}
\[
\frac12 + \frac14 + \frac18 + \frac{1}{16}
\]
की गणना क़दम दर क़दम (बाएँ से दाएँ) करो। आंशिक योगफलों का ढर्रा
देखो: और-और पद $\frac{1}{32}, \frac{1}{64}, \dots$ जोड़ते जाने पर
वे किस संख्या के पास पहुँचते हैं?
\end{exercise}

\section{समस्या: संगीत के एक ताल-खंड में छिपी भिन्नें}

\begin{problem}[{सप्ताहांत समस्या --- स्वरों के मान ऐसी भिन्नें हैं
जिनका योग ठीक-ठीक मिलना चाहिए, ताल गिनने में एक मशहूर श्रेणी छिपी
है, और बाल्कन ताल पर सम-विषम का राज चलता है}]
\label{pb:g7:fractions:1}
संगीत भिन्नों में लिखा जाता है। \emph{पूरा स्वर} $1$ ताल-खंड चलता
है (आम ``चार-चार'' ताल में); \emph{\omterm{def:g3:division:half}{आधा} स्वर} $\frac12$ चलता है;
फिर आते हैं \emph{\omterm{def:g3:division:half}{चौथाई} स्वर} $\frac14$, \emph{आठवाँ स्वर} $\frac18$
और \emph{सोलहवाँ स्वर} $\frac{1}{16}$। एक अटल नियम: \emph{एक
ताल-खंड के भीतर की सारी अवधियों का योग ठीक उस खंड के कुल के बराबर
होना चाहिए} --- चार-चार ताल में $1$। नीचे का हर सवाल इसी अध्याय का
गणित है (\cref{thm:g7:fractions:add}, \cref{met:g7:fractions:add}),
संगीत में ढला हुआ; रास्ते में तुम्हें एक हज़ार साल पहले भारतीय
विद्वानों की खोजी हुई एक गिनती-श्रेणी मिलेगी।

\medskip
\noindent\textbf{भाग I --- ऐसे स्वर जिनका योग मिलना ही चाहिए।}
\begin{enumerate}
  \item जाँचो कि इनमें से हर ताल-खंड चार-चार ताल में सही है:
        (a) \omterm{def:g3:division:half}{आधा} $+$ \omterm{def:g3:division:half}{चौथाई} $+$ \omterm{def:g3:division:half}{चौथाई}; (b) \omterm{def:g3:division:half}{चौथाई} $+$ आठवाँ $+$
        आठवाँ $+$ \omterm{def:g3:division:half}{आधा}।
  \item एक नक़लनवीस ने लिखा: \omterm{def:g3:division:half}{आधा} $+$ \omterm{def:g3:division:half}{चौथाई} $+$ आठवाँ। ताल-खंड का
        कितना भाग लिखा गया, और कौन-सा एक स्वर उसे पूरा करेगा?
  \item एक पूरे ताल-खंड में कितने सोलहवें स्वर आते हैं? एक आठवें
        स्वर में कितने सोलहवें?
  \item किसी स्वर के बाद लगा \emph{बिंदु} उसमें उसका \omterm{def:g3:division:half}{आधा} मान और
        जोड़ देता है: बिंदुदार \omterm{def:g3:division:half}{आधा} $\frac12 + \frac14$ है, और
        बिंदुदार \omterm{def:g3:division:half}{चौथाई} $\frac14 + \frac18$। दोनों मान निकालो।
  \item वाल्ट्ज़ तीन-चार ताल में लिखा जाता है: हर ताल-खंड का कुल
        $\frac34$ होना चाहिए। जाँचो कि अकेला बिंदुदार \omterm{def:g3:division:half}{आधा} एक
        वाल्ट्ज़ ताल-खंड भर देता है, और सिर्फ़ \omterm{def:g3:division:half}{चौथाई} तथा आठवें
        स्वरों से बने दो और सही वाल्ट्ज़ ताल-खंड लिखो।
\end{enumerate}

\medskip
\noindent\textbf{भाग II --- नक़लनवीस की कार्यशाला।}
\begin{enumerate}[resume]
  \item एक चार-चार ताल-खंड में है: बिंदुदार \omterm{def:g3:division:half}{चौथाई} $+$ आठवाँ $+$
        \omterm{def:g3:division:half}{चौथाई}। क्या कम पड़ रहा है?
  \item कौन ज़्यादा देर चलता है, बिंदुदार आठवाँ या \omterm{def:g3:division:half}{चौथाई} स्वर?
        सोलहवें स्वरों में गिनकर तुलना करो।
  \item \emph{जोड़-रेखा} कई अवधियों को आपस में चिपकाकर एक लंबी
        आवाज़ बना देती है: आठवें से जुड़ा हुआ \omterm{def:g3:division:half}{आधा} स्वर कितनी देर
        चलता है? वह कितने सोलहवें हुए?
  \item एक ताल-खंड के स्वरों का योग $\frac{11}{16}$ है; बाक़ी
        सन्नाटा है (एक \emph{विराम})। विराम कितनी देर का है?
  \item एक ढोलकिया एक \omterm{def:g3:division:half}{चौथाई} स्वर (ताल-खंड का $\frac{4}{16}$)
        सिर्फ़ आठवें ($\frac{2}{16}$) और सोलहवें ($\frac{1}{16}$)
        स्वरों से भरता है, और क्रम मायने रखता है: जैसे
        आठवाँ--सोलहवाँ--सोलहवाँ, या सोलहवाँ--आठवाँ--सोलहवाँ।
        जाँचो कि दोनों उदाहरण सही हैं, फिर \emph{सारी} संभव लयें
        गिनाओ। वे कितनी हैं?
\end{enumerate}

\medskip
\noindent\textbf{भाग III --- एक हज़ार साल पुरानी श्रेणी, और एक
विषम ताल।}
\begin{enumerate}[resume]
  \item आओ लंबे स्वरों के लिए ढोलकिये की लयें गिनें। हर लय या तो
        सोलहवें पर ख़त्म होती है या आठवें पर; समझाओ कि इससे यह
        नियम क्यों मिलता है: ($n$ सोलहवें भरने वाली लयें) $=$
        ($n - 1$ भरने वाली लयें) $+$ ($n - 2$ भरने वाली लयें)। एक
        सोलहवें के लिए $1$ लय और दो के लिए $2$ से शुरू करके
        $n = 8$ तक सारणी बनाओ: एक आधे स्वर को कितनी लयों में
        बजाया जा सकता है?
  \item तुम्हें मिली संख्याएँ --- $1, 2, 3, 5, 8, 13, 21, 34$ ---
        भारतीय विद्वान हेमचंद्र ने सन 1150 के आसपास ठीक ऐसी ही
        लयें गिनते हुए छापी थीं, इटली में फ़िबोनाच्ची के उन्हें
        लिखने से दो पीढ़ी पहले। उनके बनने का नियम एक वाक्य में
        लिखो, और अपनी सारणी की आख़िरी तीन प्रविष्टियों पर उसे
        जाँचो।
  \item ``स्विंग'' शैली दो बराबर आठवें स्वरों की जगह एक
        लंबा--छोटा जोड़ा रख देती है: पहले $\frac16$, फिर
        $\frac{1}{12}$। उभयनिष्ठ हर $12$ लेकर जाँचो कि यह
        अदला-बदली सही है (कुल अवधि वही रहती है)।
  \item \emph{त्रिक} तीन बराबर स्वरों को एक \omterm{def:g3:division:half}{चौथाई} स्वर में ठूँस
        देता है। हर स्वर कितनी देर का हुआ (अपने उत्तर को $3$ से
        गुणा करके जाँचो, \cref{prop:g7:fractions:mult})? वही सवाल
        एक आधे स्वर को भरने वाले तीन बराबर स्वरों के लिए।
  \item बाल्कन नृत्य-संगीत सात-आठ ताल इस्तेमाल करता है: हर
        ताल-खंड का कुल $\frac78$ होता है। जाँचो कि बिंदुदार \omterm{def:g3:division:half}{चौथाई}
        $+$ \omterm{def:g3:division:half}{चौथाई} $+$ \omterm{def:g3:division:half}{चौथाई} इसे भर देता है। फिर समझाओ कि सिर्फ़
        आधे और \omterm{def:g3:division:half}{चौथाई} स्वरों का \emph{कोई भी} जोड़, चाहे कितने भी
        हों, सात-आठ ताल-खंड को कभी नहीं भर सकता। (आठवें स्वरों
        में गिनो और सम तथा विषम के बारे में सोचो।)

\end{enumerate}
\end{problem}
